\subsubsection{Описание функции валидации параметров экзопланеты}

Для обеспечения корректности вводимых параметров экзопланеты в разрабатываемой системе реализована функция \texttt{validatePlanetParameters}. Данная функция является критически важным компонентом системы и написана на языке JavaScript. На рисунке \ref{fig:validation_code} представлена программная реализация данной функции.

Функция принимает три ключевых параметра:
\begin{itemize}
    \item \texttt{radius} --- радиус экзопланеты в единицах радиуса Земли (R⊕);
    \item \texttt{temperature} --- температура поверхности экзопланеты в Кельвинах (K);
    \item \texttt{density} --- плотность экзопланеты относительно плотности Земли.
\end{itemize}

Внутри функции создаётся объект \texttt{validation}, содержащий два поля:
\begin{itemize}
    \item \texttt{isValid: true} --- флаг, указывающий на валидность всех параметров;
    \item \texttt{errors: []} --- массив для хранения сообщений об ошибках.
\end{itemize}

Как видно из рисунка \ref{fig:validation_code}, в функции выполняется последовательная проверка каждого параметра:

\begin{enumerate}
    \item \textbf{Проверка радиуса:}
    \begin{itemize}
        \item допустимый диапазон: от 0.1 до 3.0 R⊕;
        \item при выходе за пределы диапазона устанавливается \texttt{isValid = false};
        \item в массив ошибок добавляется соответствующее сообщение.
    \end{itemize}

    \item \textbf{Проверка температуры:}
    \begin{itemize}
        \item допустимый диапазон: от 50K до 1000K;
        \item значения вне диапазона приводят к установке \texttt{isValid = false};
        \item формируется информативное сообщение об ошибке.
    \end{itemize}

    \item \textbf{Проверка плотности:}
    \begin{itemize}
        \item допустимый диапазон: от 0.01 до 3.0 относительно плотности Земли;
        \item некорректные значения отмечаются как ошибка;
        \item добавляется соответствующее пояснение в массив ошибок.
    \end{itemize}
\end{enumerate}

Выбранные диапазоны значений основаны на актуальных астрофизических данных:
\begin{itemize}
    \item диапазон радиусов охватывает планеты от сверхземель до мини-нептунов;
    \item температурный диапазон включает зоны потенциальной обитаемости различных звёздных систем;
    \item диапазон плотностей учитывает различные типы планет: от газовых гигантов до каменистых планет.
\end{itemize}

Как показано на рисунке \ref{fig:validation_code}, функция реализует принцип раннего обнаружения ошибок, что позволяет:
\begin{itemize}
    \item предотвратить выполнение расчётов с некорректными данными;
    \item обеспечить информативную обратную связь пользователю;
    \item поддерживать целостность данных в системе.
\end{itemize}

После выполнения всех проверок функция возвращает объект \texttt{validation}, содержащий результат валидации и список возможных ошибок. Это позволяет вызывающему коду принять решение о дальнейших действиях на основе полученной информации.

\begin{figure}[H]
    \centering
    \includegraphics[width=0.8\textwidth]{image/code/validation_code.png}
    \caption{Реализация функции валидации параметров экзопланеты}
    \label{fig:validation_code}
\end{figure}

\subsubsection{Математическая формула}

ESI рассчитывается как произведение компонентных индексов подобия:

\begin{equation}
\label{eq:esi_formula}
ESI = \prod_{i=1}^{n} (1 - |\frac{x_i - x_{0i}}{x_i + x_{0i}}|)^{w_i/n}
\end{equation}

где:
\begin{itemize}
    \item $x_i$ --- значение параметра экзопланеты;
    \item $x_{0i}$ --- эталонное значение параметра (земное);
    \item $w_i$ --- весовой коэффициент параметра;
    \item $n$ --- количество учитываемых параметров.
\end{itemize}

\subsubsection{Входные параметры}

Функция получает на вход объект params, содержащий основные параметры экзопланеты: radius --- радиус планеты в радиусах Земли и temperature --- температуру поверхности в Кельвинах. Для сравнения с эталонными земными значениями используется объект referenceParams, где radius = 1.0, а temperature = 288K. Весовые коэффициенты параметров, определяющие их значимость при расчетах, задаются в объекте weights следующим образом: для radius коэффициент составляет 0.57, а для temperature --- 0.43.

\subsubsection{Диапазон значений}

Результирующий индекс подобия Земле (ESI) находится в интервале от 0 до 1. Значение 1.0 соответствует полному совпадению с параметрами Земли, а 0.0 указывает на полное несоответствие земным параметрам. Интерпретация различных значений ESI приведена в таблице \ref{tab:esi_classification}.

\begin{table}[h]
\captionsetup{justification=raggedright,singlelinecheck=false}
\caption{Классификация экзопланет по индексу подобия Земле}
\label{tab:esi_classification}
\begin{tabular}{|l|l|}
\hline
\textbf{Диапазон ESI} & \textbf{Интерпретация} \\
\hline
0.9 - 1.0 & Очень высокое подобие Земле \\
0.8 - 0.9 & Высокое подобие Земле \\
0.6 - 0.8 & Среднее подобие Земле \\
0.4 - 0.6 & Низкое подобие Земле \\
0.0 - 0.4 & Очень низкое подобие Земле \\
\hline
\end{tabular}
\end{table}

\subsubsection{Примечания}

\begin{enumerate}
    \item Функция использует нормализованные значения параметров для обеспечения сопоставимости различных физических характеристик.
    \item Весовые коэффициенты отражают относительную важность каждого параметра в определении пригодности планеты для жизни.
    \item Формула учитывает как абсолютные различия между параметрами, так и их относительные значения.
\end{enumerate}

\subsubsection{Анализ факторов обитаемости экзопланеты}

Для комплексной оценки потенциальной обитаемости экзопланеты в системе реализована функция analyzeHabitabilityFactors. Данная функция входит в состав программного модуля анализа экзопланетных характеристик и предназначена для определения пригодности планеты к потенциальному существованию на ней жизни, схожей с земной.

Функция принимает объект параметров экзопланеты и анализирует три ключевых фактора: температуру, гравитацию и атмосферу. Каждый фактор оценивается как количественно (по шкале от 0 до 1), так и качественно (формируется текстовое описание).

При оценке температурного режима функция проверяет, находится ли температура поверхности планеты в диапазоне, подходящем для существования воды в жидком состоянии (от 273K до 373K). Если температура соответствует этому условию, вычисляется оценка, основанная на близости к оптимальной температуре 288K (средняя температура поверхности Земли). Для планет с температурой в указанном диапазоне формируется описание "Температура в диапазоне жидкой воды". Если же температура выходит за границы интервала, оценка значительно снижается, а текстовое описание указывает на неблагоприятные условия для жизни.

Оценка гравитационных условий осуществляется на основе радиуса планеты. Функция вычисляет относительную силу тяжести и сравнивает ее с земными условиями. Если гравитация экзопланеты находится в пределах от 0.8 до 1.2 земной, формируется описание "Гравитация близка к земной". При существенных отклонениях от этого диапазона создается соответствующая характеристика, указывающая на значительные отличия от земных условий.

Третий фактор – атмосфера – оценивается с использованием параметра плотности. В зависимости от значения плотности атмосферы относительно земной формируется либо заключение "Присутствует существенная атмосфера" (при плотности более 3.0), либо "Атмосфера разрежена или отсутствует" (при меньших значениях). Числовая оценка рассчитывается с учетом пороговых значений, определенных на основе планетологических исследований.

По результатам анализа функция возвращает объект, содержащий комплексную информацию о трех рассмотренных факторах. Каждый фактор представлен числовой оценкой и текстовым описанием. Такая форма представления результатов позволяет как автоматизировать дальнейшую обработку данных, так и обеспечить понятную интерпретацию для пользователей системы.

Данная функция является важным компонентом разрабатываемой системы и вносит существенный вклад в общий процесс анализа экзопланет, предоставляя структурированную оценку их потенциальной обитаемости на основе доступных астрофизических данных.

\subsubsection{Расчет вероятности наличия воды}

В рамках системы анализа экзопланет реализована функция calculateWaterProbability, предназначенная для оценки вероятности наличия воды в жидком состоянии на поверхности планеты. Функция учитывает два ключевых параметра: температуру поверхности и радиус планеты.

В основе алгоритма лежит анализ температурного диапазона, благоприятного для существования жидкой воды. В условиях земного атмосферного давления этот диапазон составляет от 273K (0°C) до 373K (100°C). Если температура экзопланеты находится в данных пределах, вычисляется базовая вероятность наличия воды. Максимальное значение этой вероятности соответствует середине температурного диапазона и снижается по мере приближения к его границам. При температурах, выходящих за указанные рамки, базовая вероятность устанавливается равной нулю.

После определения базовой вероятности на основе температуры выполняется корректировка с учетом радиуса планеты, который косвенно характеризует ее гравитационные свойства. Данная корректировка отражает тот факт, что планеты с радиусом, существенно отличающимся от земного, могут иметь атмосферу с давлением, при котором диапазон температур существования жидкой воды отличается от принятого за стандарт. Так, на планетах с низкой гравитацией (малый радиус) вода может испаряться при более низких температурах из-за пониженного атмосферного давления, а на планетах с высокой гравитацией (большой радиус) точка кипения воды может быть выше.

Результат работы функции ограничивается значениями от 0 до 1, где 0 означает полное отсутствие вероятности наличия жидкой воды, а 1 – максимальную вероятность. Это позволяет использовать данную оценку как в составе комплексного анализа обитаемости экзопланеты, так и в качестве самостоятельного показателя при исследовании потенциально пригодных для жизни миров.

\subsubsection{Результаты тестирования}

В ходе тестирования были проверены все основные функции системы и подтверждена их корректная работа. Особое внимание было уделено точности математических расчетов и качеству визуализации данных. Тестирование показало, что система:

\begin{itemize}
    \item корректно валидирует входные параметры экзопланет;
    \item точно выполняет расчеты индекса подобия Земле и других показателей;
    \item качественно визуализирует результаты анализа;
    \item обеспечивает удобный и информативный интерфейс.
\end{itemize}

Все выявленные в процессе тестирования недочеты были успешно устранены, что позволяет говорить о готовности системы к практическому использованию в области анализа потенциально обитаемых экзопланет.

\begin{table}[H]
\captionsetup{justification=raggedright,singlelinecheck=false}
\caption{Среднее время отклика на основные действия в системе анализа экзопланет}
\label{table_2}
\begin{tabular}{|p{9cm}|r|}
    \hline
    \textbf{Функциональность} & \textbf{Среднее время отклика, сек} \\
    \hline
    Валидация параметров экзопланеты & 0,45 \\
    \hline
    Расчет индекса подобия Земле (ESI) & 0,82 \\
    \hline
    Построение 3D-модели экзопланеты & 1,35 \\
    \hline
    Генерация графиков и диаграмм & 1,12 \\
    \hline
    Поиск похожих экзопланет в базе & 0,98 \\
    \hline
    Расчет вероятности наличия воды & 0,89 \\
    \hline
    Анализ атмосферных параметров & 0,93 \\
    \hline
    \textbf{Среднее значение по всем операциям} & \textbf{0,93} \\
    \hline
\end{tabular}
\end{table}

Исходя из результатов тестирования в таблице \ref{table_2}, можно сделать вывод, что время отклика на все виды запросов находится в пределах нормы.

Для оценки ресурсных показателей системы анализа экзопланет использовались встроенные инструменты браузера Google Chrome --- в частности, Chrome DevTools Performance Profiler. Он позволяет в реальном времени отслеживать использование ресурсов при взаимодействии пользователя с интерфейсом: от ввода параметров до финального отображения результатов анализа.

Оптимизация производительности клиентской части системы предоставляет ряд важных преимуществ, включая снижение времени расчетов, повышение отзывчивости интерфейса и комфортность использования приложения на различных устройствах.

Performance Profiler в Chrome позволяет:
\begin{itemize}
    \item визуализировать процесс расчета параметров обитаемости;
    \item определить этапы построения 3D-моделей и графиков;
    \item отслеживать время отклика на изменение параметров;
    \item измерить метрики, характеризующие восприятие производительности.
\end{itemize}

Информация, собираемая при профилировании:
\begin{itemize}
    \item LCP (Largest Contentful Paint) --- фиксирует время отображения 3D-модели и основных графиков;
    \item CLS (Cumulative Layout Shift) --- показывает стабильность интерфейса при обновлении данных;
    \item INP (Interaction to Next Paint) --- измеряет задержку при интерактивном взаимодействии с 3D-моделью.
\end{itemize}

После проведения анализа производительности интерфейса системы были получены следующие показатели:

\begin{itemize}
    \item LCP: 1,35 секунды --- время построения 3D-модели экзопланеты;
    \item CLS: 0,02 --- минимальные смещения при обновлении графиков;
    \item INP: 75 мс --- быстрый отклик на вращение 3D-модели.
\end{itemize}

Для проверки надёжности системы анализа экзопланет была применена упрощённая оценка надёжности, основанная на количественной фиксации сбоев при многократном использовании системы.

Для оценки надёжности программного средства использовалась следующая формула:

\begin{equation}
\label{eq:reliability}
P = 1 - \frac{n}{N},
\end{equation}

где $n$ --- количество зафиксированных сбоев,  
$N$ --- общее число запусков/испытаний системы.

В ходе функционального тестирования было выполнено 100 последовательных запусков ключевых сценариев:
\begin{itemize}
    \item ввод параметров экзопланеты;
    \item расчет индекса подобия Земле;
    \item построение 3D-модели;
    \item генерация графиков;
    \item поиск похожих экзопланет.
\end{itemize}

В процессе тестирования было зафиксировано 3 нестабильных поведения:
\begin{itemize}
    \item один случай некорректного построения 3D-модели;
    \item один случай задержки при расчете индекса ESI;
    \item один случай ошибки при поиске похожих экзопланет.
\end{itemize}

Тогда, согласно формуле:

\begin{equation}
\label{eq:reliability-result}
P = 1 - \frac{3}{100} = 0{,}97.
\end{equation}

Полученное значение $0{,}97$ свидетельствует о высокой степени надёжности системы анализа экзопланет. Это означает, что система корректно функционирует в 97\% случаев при стандартных пользовательских сценариях.

Основные аспекты надёжности системы:
\begin{itemize}
    \item точность расчета параметров обитаемости;
    \item стабильность визуализации данных;
    \item корректность поиска похожих экзопланет;
    \item устойчивость работы 3D-рендеринга.
\end{itemize}

\subsection{Вывод}

В ходе разработки интеллектуальной системы анализа обитаемости экзопланет были использованы современные веб-технологии, включая Three.js для 3D-визуализации и Plotly для построения графиков. Были проведены работы по разработке математических моделей, алгоритмов анализа и визуального представления результатов. Основной целью проекта было создание инструмента, позволяющего исследователям эффективно оценивать потенциальную обитаемость экзопланет на основе их физических параметров.

В процессе разработки системы были учтены требования к точности расчетов, надежности и удобству использования. Были реализованы основные функциональные возможности, включая:
\begin{itemize}
    \item валидацию входных параметров экзопланет;
    \item расчет индекса подобия Земле (ESI);
    \item анализ вероятности наличия жидкой воды;
    \item построение интерактивных 3D-моделей;
    \item генерацию наглядных графиков и диаграмм;
    \item поиск и сравнение похожих экзопланет.
\end{itemize} 